\chapter{使用说明}

阅读时可将示例与编辑器并排放置：先按原样运行，再修改一项输入，记录结果与原因。练习的目标是能够独立复现并解释行为。
读到报错时也不用急着翻过去，那里往往藏着比成功输出更有用的线索。

\section*{开工前，把三件东西对上}

先准备 Python 3 和 PyCharm，再为本书练习建一个独立项目。
项目里的 \code{.venv} 放解释器与依赖，自己的脚本放在它外面。
本版运行验证使用 Python 3.12，具体依赖和验证范围见项目包中的
\code{RUN_EXAMPLES.md}。换解释器版本以后，先重跑小样例，再运行项目。

运行前请核对三处解释器设置：编辑器选中的项目解释器、运行配置所用的解释器、终端启动的解释器。
如图~\ref{fig:usage-interpreter}，把终端里看到的路径与项目设置对照，
通常比反复点安装按钮更快找到问题。
\begin{figure}[htbp]
  \centering
  \BookImage{interpreter}{\linewidth}
  \caption{项目解释器、运行配置与终端解释器的路径对应关系。}
  \label{fig:usage-interpreter}
\end{figure}

\section*{代码块怎样读}

带行号的 Python 代码块是脚本。文件名写在上方，复制正文时把行号留在书里，
保存为 UTF-8 编码的 \code{.py} 文件。缩进每层使用四个空格。
几行属于同一个缩进块，可以沿着左侧空白一眼看出来，别用全角空格补齐。

\code{>>>} 是交互式解释器的主提示符，跟练时只输入它右边的代码。
终端命令放在标有 terminal 的代码块里，在终端执行。
“运行核对”中的输出只用于对照，代码中注明“取消注释以复现”的语句则是刻意留下的错误入口。
先跑正常版本，再单独打开一处错误，免得几个问题挤到同一个回溯里。

注释解释这一段为什么这样写。改动代码时，注释也要跟着核对。
如果示例只展示某个项目中的函数，标题会标明它在源码中的位置，
完整的导入、入口和依赖都在对应项目目录里。

\section*{示例运行六步法}

\begin{enumerate}
  \item 找到文件名和所在目录，先看需要哪些库、样本或服务。
  \item 原样运行，对照输出。路径和耗时随机器变化，先核对含义与顺序。
  \item 用断点看关键变量，确认是哪个输入走到了哪条分支。
  \item 每次只改一个条件，例如空列表、缺失字段、负数或不存在的路径。
  \item 给失败留下输入与回溯，再按附录 C 的路径排查。
  \item 重新跑原例和变体，把修复前后的结果一起记下。
\end{enumerate}

\begin{myunderstanding}
将练习看作一次小型实验：输入是条件，代码是步骤，输出是可核对的结果。运行成功以后，再改变一项条件，解释结果为何随之变化。
\end{myunderstanding}

\section*{先留下自己的运行现场}

下面四行适合放在新练习最前面。它们记录解释器版本、位置和工作目录，
方便将终端与 PyCharm 的运行结果放在一起比较。
提交代码前，可再把这些临时诊断输出收进日志或删除。
\begin{lstlisting}[title={\code{check_environment.py}}]
import sys
from pathlib import Path

# 三个现场信息，核对解释器与相对路径的起点。
print(sys.version.split()[0])
print(sys.executable)
print(Path.cwd())
\end{lstlisting}

\section*{项目包怎样用}

主文件 \code{main.tex} 管篇章顺序，\code{preamble.tex} 管字体、配色、提示框与目录。
\code{chapters} 放正文，\code{examples} 放章节脚本，\code{projects} 放四阶段项目。
配图在 \code{figures}，生成它们的代码在 \code{tools}。
可把整个项目目录解压后再编译，避免只拿一个主文件漏掉图和输入文件。

从 \code{examples/ch11} 运行带参数的日志例子时，使用同目录的
\code{sample.log}。项目三则把输入与输出路径都写进命令。
碰到找不到文件，先打印当前工作目录，再看参数究竟指向哪里。
完整运行方法见 \code{RUN_EXAMPLES.md}。

\section*{图、提示框和边界条件}

蓝框「核心直觉」把概念落到一个能记住的形象上，黄框「PyCharm 实战」给出可操作的观察点，
红框「踩坑记录」则留下出错条件和修复路线。比喻用来建立直觉，
真正的边界以代码、输出和文中说明为准。

界面图注明“示意”时，画的是观察关系。真实 PyCharm 的按钮位置和快捷键以所选 Keymap 为准。
浏览器、桌面 GUI、外部数据库和 Docker 例子需要相应环境。
本版对这些部分写清了验证范围，连接参数都由读者的测试环境提供。

\section*{按任务安排阅读}

第一次接触 Python，可从第一篇开始，每节保留一个能独立运行的文件。
已经写过脚本的读者，可先按第 7、8 章补充测试、日志和版本记录，再查阅相关语言规则。
要处理数据，就先走文件与持久化，再接第 16 章和项目三。
要写网络任务，先把第 13 章的本机示例跑顺，再加入并发。

书里的练习可以变成自己的笔记。每章末尾，可写下三行，今天跑通了什么，
改了哪个输入，以及哪条报错以后想一眼认出来。
久而久之，熟练就从这些有来路的细节里长出来。

\paragraph{本次代码修订说明} 灰色折行箭头表示版面续行，不属于 Python 语法。五个较长示例已改成括号内合法换行或中间变量，随书 examples 中的源文件与正文同步。PDF 复制仍可能带入行号、箭头或改变缩进，请优先从随书 .py 文件运行。红底代码为故意失败或产生错误结果的反例，对应修复版本使用普通底色。
